\setcounter{section}{2}

\Needspace{5\baselineskip}
\hypertarget{ux72b6ux6001ux5206ux5e03ux4e0eux5360ux7528ux5ea6ux91cf}{%
\section{状态分布与占用度量}\label{ux72b6ux6001ux5206ux5e03ux4e0eux5360ux7528ux5ea6ux91cf}}

\Needspace{5\baselineskip}
\hypertarget{ux4e00ux72b6ux6001ux5206ux5e03ux548cux5360ux7528ux5ea6ux91cfux7684ux5b9aux4e49}{%
\subsection{一、状态分布和占用度量的定义}\label{ux4e00ux72b6ux6001ux5206ux5e03ux548cux5360ux7528ux5ea6ux91cfux7684ux5b9aux4e49}}

\Needspace{5\baselineskip}
状态分布定义为智能体处于状态 \(s\) 的折扣后概率：

\Needspace{5\baselineskip}
定义了 \(\nu^\pi(s)\)，即在策略 \(\pi\) 下的折扣状态访问分布：

\[
\nu^\pi(s) = (1-\gamma)\sum_{t=0}^{\infty}\gamma^t P_t^\pi(s)
\]

\(P_t^\pi(s)\)：表示在时刻 \(t\)，智能体处于状态 \(s\) 的概率。

\(\gamma\)：折扣因子（Discount Factor），\(0 \le \gamma < 1\)。

\(1-\gamma\)：归一化因子（Normalization Factor）。

\Needspace{5\baselineskip}
进而可以定义占用度量，即智能体在s发出动作a的折扣后概率：

\[
\rho^\pi(s,a) = (1-\gamma)\sum_{t=0}^{\infty}\gamma^t P_t^\pi(s)\pi(a|s)
\]

\Needspace{5\baselineskip}
此时有关系：

\[
\rho^\pi(s,a) = \nu^\pi(s)\pi(a|s)
\]

\Needspace{5\baselineskip}
\hypertarget{ux4e8cux539fux7406ux63a8ux5bfc}{%
\subsection{二、原理推导}\label{ux4e8cux539fux7406ux63a8ux5bfc}}

\begin{enumerate}
\def\labelenumi{\arabic{enumi}.}
\tightlist
\item
  为什么需要折扣因子？
\end{enumerate}

\Needspace{5\baselineskip}
我们通常希望最大化累积折扣回报：

\[
J(\pi) = \mathbb{E}_\pi\left[\sum_{t=0}^{\infty}\gamma^t R(s_t,a_t)\right]
\]

\Needspace{5\baselineskip}
利用期望的线性性质，我们可以把求和号拿到期望外面：

\[
J(\pi) = \sum_{t=0}^{\infty}\gamma^t\mathbb{E}_\pi[R(s_t,a_t)]
\]

\[
J(\pi) = \sum_{t=0}^{\infty}\gamma^t\sum_{s,a}P_t^\pi(s)\pi(a|s)R(s,a)
\]

\Needspace{5\baselineskip}
现在，我们将对于 \(s\) 和 \(a\) 的求和提到最前面：

\[
J(\pi) = \sum_{s,a}R(s,a)\left(\sum_{t=0}^{\infty}\gamma^t P_t^\pi(s)\pi(a|s)\right)
\]

\Needspace{5\baselineskip}
括号中的这一项是未归一化的占用度量。为了让这一项变成概率分布（和为 1），我们乘以并除以 \((1-\gamma)\)：

\[
J(\pi) = \frac{1}{1-\gamma}\sum_{s,a}R(s,a)\left((1-\gamma)\sum_{t=0}^{\infty}\gamma^t P_t^\pi(s)\pi(a|s)\right)
\]

\Needspace{5\baselineskip}
其中括号中的部分就是 \(\rho^\pi(s,a)\)，因此：

\[
J(\pi) = \frac{1}{1-\gamma}\sum_{s,a}\rho^\pi(s,a)R(s,a) = \frac{1}{1-\gamma}\mathbb{E}_{(s,a)\sim\rho^\pi}[R(s,a)]
\]

在强化学习中，我们的目标往往是找到一个好的策略。策略可以用概率分布表示，优化目标正是让回报 \(R\) 在 \(\rho^\pi\) 的概率分布下的数学期望最大化。这样一来，问题就转化为寻找这样的概率分布；状态-动作对在分布中的概率，等于该状态-动作对对总价值的贡献权重，符合我们对马尔可夫决策过程的认知。

\begin{enumerate}
\def\labelenumi{\arabic{enumi}.}
\setcounter{enumi}{1}
\tightlist
\item
  为什么需要归一化因子？
\end{enumerate}

这是一个非常关键的数学细节。我们需要证明 \(\nu^\pi(s)\) 是一个合法的概率分布，即其对所有状态求和必须为 1。

\Needspace{5\baselineskip}
推导如下：我们对所有状态 \(s\) 求和：

\[
\sum_s \nu^\pi(s) = \sum_s\left[(1-\gamma)\sum_{t=0}^{\infty}\gamma^t P_t^\pi(s)\right]
\]

\Needspace{5\baselineskip}
交换求和顺序（假设级数收敛）：

\[
= (1-\gamma)\sum_{t=0}^{\infty}\gamma^t\left(\sum_s P_t^\pi(s)\right)
\]

\Needspace{5\baselineskip}
因为无论 \(t\) 是多少，智能体一定会在某个状态，所以 \(\sum_s P_t^\pi(s)=1\)。公式变为：

\[
= (1-\gamma)\sum_{t=0}^{\infty}\gamma^t \cdot 1 = 1
\]

\Needspace{5\baselineskip}
\hypertarget{ux4e09ux5360ux7528ux5ea6ux91cfux7684ux6570ux5b66ux542bux4e49ux6309ux51e0ux4f55ux5206ux5e03ux5728ux65f6ux95f4ux4e0aux91c7ux6837}{%
\subsection{三、占用度量的数学含义：按几何分布在时间上采样}\label{ux4e09ux5360ux7528ux5ea6ux91cfux7684ux6570ux5b66ux542bux4e49ux6309ux51e0ux4f55ux5206ux5e03ux5728ux65f6ux95f4ux4e0aux91c7ux6837}}

通常我们理解的``访问概率''，往往指的是长期来看，在所有时间步中，我有百分之多少的时间待在这个状态（即简单的频率统计）。

但加上折扣因子 \(\gamma\) 后，现在的 \(\rho^\pi(s,a)\) 显然不再是简单的``时间频率''了（因为它重前期、轻后期）。

之所以依然称之为``概率''，是因为我们改变了``采样时间 \(t\)''的概率分布。

传统视角（无折扣，\(\gamma=1\)）

你假设每一个时刻 \(t\) 被抽到的概率是相等的（这在无限时间内其实无法实现数学上的均匀分布，通常取极限）。这就导致了我们常说的``平稳分布''（Stationary Distribution），即单纯的时间平均占比。

折扣视角（有折扣，\(\gamma<1\)）

这里的 \(\rho^\pi(s,a)\) 实际上定义了一个新的随机实验。在这个实验中，我们不是随机均匀地抽取一个时刻，而是按照几何分布（Geometric Distribution）来抽取时刻 \(t\)。

\Needspace{5\baselineskip}
采样规则：

\begin{itemize}
\tightlist
\item
  抽到 \(t=0\) 的概率是 \((1-\gamma)\)。
\item
  抽到 \(t=1\) 的概率是 \((1-\gamma)\gamma\)。
\item
  抽到 \(t=2\) 的概率是 \((1-\gamma)\gamma^2\)。
\item
  \ldots{}
\item
  抽到时刻 \(t\) 的概率是 \(P(T=t)=(1-\gamma)\gamma^t\)，其中 \(T\) 表示采样时刻。
\end{itemize}

\Needspace{5\baselineskip}
现在的物理含义是：

如果不确定在哪个时间点观察，而是按照上述``越近越重要''的规则随机抽取一个瞬间，那么在这个瞬间，智能体位于 \((s,a)\) 的概率是多少？

\Needspace{5\baselineskip}
也可以这样理解：

\Needspace{5\baselineskip}
为了更直观地理解，我们可以引入一个随机终止的设定（这也是 \(\gamma\) 最经典的物理意义）：

\Needspace{5\baselineskip}
假设这个游戏不是无限进行的，而是在每一步结束后，上帝会掷一个骰子：

\begin{itemize}
\tightlist
\item
  有 \(\gamma\) 的概率让游戏继续。
\item
  有 \(1-\gamma\) 的概率让游戏突然结束。
\end{itemize}

\Needspace{5\baselineskip}
那么，\(\rho^\pi(s,a)\) 衡量的就是：

当游戏``突然结束''的那一瞬间（Game Over Snapshot），智能体正处于状态 \(s\) 且正在做动作 \(a\) 的概率。

\Needspace{5\baselineskip}
这个解释完美契合了公式：

\begin{itemize}
\tightlist
\item
  游戏恰好在时刻 \(t\) 结束的概率是 \((1-\gamma)\gamma^t\)。
\item
  在时刻 \(t\) 智能体位于 \((s,a)\) 的概率是 \(P_t^\pi(s,a)\)。
\item
  把所有可能结束的时刻加起来，就是终止时看到 \((s,a)\) 的总概率。
\end{itemize}

因为游戏一定会结束（在某一个时刻），所以所有 \((s,a)\) 的这种概率加起来一定等于 1。因此，它是一个完美的概率分布。

\Needspace{5\baselineskip}
\hypertarget{ux56dbux72b6ux6001ux5206ux5e03ux7684ux9012ux5f52ux8868ux793a}{%
\subsection{四、状态分布的递归表示}\label{ux56dbux72b6ux6001ux5206ux5e03ux7684ux9012ux5f52ux8868ux793a}}

当我们计算 \(\nu^\pi(s')\) 时，我们统计的是从 \(t=0\) 到 \(t=\infty\) 的所有时间步中，智能体出现在 \(s'\) 的加权总和。

\Needspace{5\baselineskip}
对于任何一次出现，它发生的时间只有两种可能：

\begin{enumerate}
\def\labelenumi{\arabic{enumi}.}
\tightlist
\item
  它发生在 \(t=0\)：这就是``初始状态''。
\item
  它发生在 \(t>0\)：这就是``从上一时刻转移过来的''。
\end{enumerate}

在 \(s'\) 的总概率 = 作为起点的权重 + 从别处转移来的权重。

\[
\nu^\pi(s') = (1-\gamma)\nu_0(s') + \gamma\int P(s'|s,a)\pi(a|s)\nu^\pi(s)\,ds\,da
\]

第一项 \((1-\gamma)\nu_0(s')\) 是初始贡献，第二项 \(\gamma\int P(s'|s,a)\pi(a|s)\nu^\pi(s)\,ds\,da\) 是转移贡献。

其中 \(\nu_0(s')\) 指的是初始状态下在 \(s'\) 的概率 \(P(s_0=s')\)，和 \(\nu^\pi\) 不是同一个函数。

第一项：\((1-\gamma)\nu_0(s')\)

\begin{itemize}
\tightlist
\item
  含义：智能体在 \(t=0\) 时刻就位于 \(s'\) 的加权概率。
\item
  \(\nu_0(s')\)：初始状态分布，即 \(P(s_0=s')\)。
\item
  \((1-\gamma)\)：这是一个时间权重的归一化系数。
\item
  回忆 \(\nu^\pi\) 的定义：\(\nu^\pi(s)=(1-\gamma)\sum_{t=0}^{\infty}\gamma^t P(s_t=s)\)。
\item
  当 \(t=0\) 时，\(\gamma^0=1\)。
\item
  所以 \(t=0\) 这一时刻对总分布的贡献就是系数 \((1-\gamma)\) 乘以概率 \(\nu_0(s')\)。
\end{itemize}

第二项：\(\gamma\int P(s'|s,a)\pi(a|s)\nu^\pi(s)\,ds\,da\)

\begin{itemize}
\tightlist
\item
  含义：智能体从上一时刻（任意状态 \(s\)）经过一步转移到达 \(s'\) 的加权概率。
\item
  \(\nu^\pi(s)\)：上一时刻在状态 \(s\) 的概率密度。
\item
  \(\pi(a|s)\)：在 \(s\) 选择动作 \(a\) 的概率。
\item
  \(P(s'|s,a)\)：执行 \(a\) 后从 \(s\) 跳到 \(s'\) 的概率。
\item
  \(\int \cdots dsda\)：对所有可能的``前任''状态和动作进行积分（如果是离散空间就是求和 \(\sum\)），把所有流向 \(s'\) 的可能性汇总。
\item
  \(\gamma\)：关键的折扣因子。
\item
  因为从 \(s\) 到 \(s'\) 经过了 1 个时间步。
\item
  如果在上一时刻（比如 \(t\)）的权重是 \(\gamma^t\)，那么到达 \(s'\) 时已经是时刻 \(t+1\)，权重变成了 \(\gamma^{t+1}\)。
\item
  相比于上一时刻，权重衰减了 \(\gamma\) 倍。所以必须乘上 \(\gamma\)。
\end{itemize}

\Needspace{5\baselineskip}
\hypertarget{ux4e94ux5360ux7528ux5ea6ux91cfux76f8ux5173ux5b9aux7406}{%
\subsection{五、占用度量相关定理}\label{ux4e94ux5360ux7528ux5ea6ux91cfux76f8ux5173ux5b9aux7406}}

进一步得出如下两个定理。

\Needspace{5\baselineskip}
定理 1：智能体分别以策略 \(\pi_1\) 和 \(\pi_2\) 与同一个 MDP 交互，得到的占用度量 \(\rho^{\pi_1}\) 和 \(\rho^{\pi_2}\) 满足：

\[
\rho^{\pi_1} = \rho^{\pi_2} \Longleftrightarrow \pi_1 = \pi_2
\]

\Needspace{5\baselineskip}
定理 2：给定一个合法占用度量 \(\rho\)，可生成该占用度量的唯一策略是：

\[
\pi_\rho = \frac{\rho(s,a)}{\sum_{a'}\rho(s,a')}
\]

\FloatBarrier
\Needspace{5\baselineskip}
\hypertarget{ux53c2ux8003ux6587ux732e-44}{%
\subsection{参考文献}\label{ux53c2ux8003ux6587ux732e-44}}

\vspace{0.25em}
\begin{itemize}
\tightlist
\item
  \href{https://hrl.boyuai.com/}{《动手学强化学习》}.
\end{itemize}

\clearpage
